\chapter{Theoretical Framework}
\label{Chapter3}
\lhead{Chapter 3. \emph{Theoretical Framework}}

The Tier~3 detector runs a small mathematical pipeline every detection
cycle: a bigram sketch becomes a 74-dimensional signal, the signal is
projected into an edge-local coupling subspace by a CCA-learned restriction
map, Mahalanobis energy scores the edge, and a global Rayleigh quotient
over a cellular sheaf Laplacian scores the whole container graph. A guarded
exponential moving average keeps the global threshold honest under normal
drift without letting a slow-drip attack poison it.

We state the classical pieces as given and cite primary
sources~\cite{hotelling1936,mahalanobis1936,hansen2019spectral,curry2014sheaves,robinson2014topological,chung1997spectral}.
No new mathematics about these objects is claimed. The contribution of this
chapter is the assembly into a deployable detector and the four theorems in
Section~\ref{sec:theorems}, which capture why that assembly is hard to
evade.

%─────────────────────────────────────────────────────────────────────────────
\section{Notation}
\label{sec:notation}
%─────────────────────────────────────────────────────────────────────────────

$C = \{c_1, \ldots, c_n\}$ is the set of monitored containers, each with
cgroup identifier $g_i \in \mathbb{N}$. A detection cycle is a window of
length $T$ indexed by $t$. At cycle $t$ each container produces a signal
$\mathbf{x}_i^{(t)} \in \mathbb{R}^d$ with $d = 74$. The communication graph
$G = (V, E)$ has $V = C$ and $E$ the set of observed inter-container edges;
we write $e = (u, v, \ell)$ with $\ell \in \{0,1,2\}$ encoding the temporal
lag. After calibration every edge carries restriction maps
$F_u^{(e)}, F_v^{(e)} \in \mathbb{R}^{k \times d}$ with $k = 15$, a
regularised covariance $\Sigma_e \in \mathbb{R}^{k \times k}$, and a
Mahalanobis threshold $\tau_e$. The global Rayleigh threshold is
$\tau_{\text{global}}$.

%─────────────────────────────────────────────────────────────────────────────
\section{Bigram sketch and the 74-dimensional signal}
\label{sec:signal}
%─────────────────────────────────────────────────────────────────────────────

Let $\Sigma$ be the set of top-25 syscall opcodes observed across the
testbed during calibration (full list in Appendix~\ref{AppendixB}). A
\emph{bigram} is an ordered pair $(s_i, s_j) \in \Sigma \times \Sigma$
recording that opcode $s_j$ immediately followed $s_i$ in the same cgroup.
In-kernel counting uses a Count-Min
Sketch~\cite{cormode2005cms} with $r{=}4$ rows and $w{=}128$ columns; the
estimate of the bigram count is $\min_\rho \mathrm{CMS}[\rho][h_\rho(\cdot)
\bmod w]$. The sketch resets every $T$~seconds, so each cycle observes
only current traffic.

The reconstructed 625-dimensional bigram vector $\mathbf{b}^{(t)}$ is
compressed into $\mathbf{x}^{(t)} = \phi(\mathbf{b}^{(t)}) \in \mathbb{R}^{74}$
through four blocks:

\begin{itemize}
  \item \textbf{Shannon entropy} (3 dims): marginal entropy of first
        opcodes, marginal entropy of second opcodes, joint bigram entropy.
  \item \textbf{PCA} (50 dims): projection onto the top-50 principal
        components of the calibration bigram set, learned once and
        frozen.
  \item \textbf{Top-$k$ marginals} (20 dims): marginal frequencies of the
        20 most common single syscalls.
  \item \textbf{Rate} (1 dim): $\log_{10}$ of total syscalls in the
        window.
\end{itemize}

The signal is then whitened per container,
$\tilde{\mathbf{x}}_i^{(t)} = (\mathbf{x}_i^{(t)} - \boldsymbol{\mu}_i)
\oslash (\mathbf{s}_i + \epsilon)$, with $\epsilon = 10^{-6}$. Diagonal
whitening is deliberate: a full $74\!\times\!74$ covariance inversion would
be rank-bounded by the per-container sample count during calibration.

\paragraph{Intuition for the shape.} The entropy block captures
\emph{disorder} in the syscall stream. The PCA block captures
\emph{structural pattern} in the bigram sketch. The marginals capture
\emph{what} the container is doing. The rate captures \emph{how much} it
is doing. A reverse shell perturbs all four blocks; a fork bomb perturbs
the rate and entropy blocks; a cross-container pivot perturbs primarily
the PCA block. Separating these blocks is what lets the sheaf Laplacian
find an attack that would be invisible to any single dimension.

%─────────────────────────────────────────────────────────────────────────────
\section{CCA restriction maps per edge}
\label{sec:cca}
%─────────────────────────────────────────────────────────────────────────────

Canonical-correlation analysis~\cite{hotelling1936} finds, for aligned
matrices $X \in \mathbb{R}^{T \times d}$ and $Y \in \mathbb{R}^{T \times d}$,
linear projections $F_X, F_Y \in \mathbb{R}^{k \times d}$ that maximise the
correlation between $F_X X^\top$ and $F_Y Y^\top$:
\begin{equation}
  (F_X^\star, F_Y^\star) = \operatorname*{arg\,max}_{F_X, F_Y}
  \frac{\operatorname{tr}\!\bigl(F_X X^\top Y F_Y^\top\bigr)}
       {\sqrt{\operatorname{tr}(F_X X^\top X F_X^\top)\,\operatorname{tr}(F_Y Y^\top Y F_Y^\top)}}.
  \label{eq:cca}
\end{equation}
The solution is a generalised eigendecomposition of the joint covariance
matrix; closed forms are in the primary source. For every observed edge
$e = (u, v, \ell)$ we form aligned calibration matrices
$X_u^{(\ell)} = [\tilde{\mathbf{x}}_u^{(1)}, \ldots,
\tilde{\mathbf{x}}_u^{(T-\ell)}]^\top$ and
$X_v^{(\ell)} = [\tilde{\mathbf{x}}_v^{(1+\ell)}, \ldots,
\tilde{\mathbf{x}}_v^{(T)}]^\top$, and solve~\eqref{eq:cca} to obtain
$F_u^{(e)}, F_v^{(e)}$. Three lags per edge are calibrated.

\paragraph{Why $k = 15$ and not $k = 50$.}  The CCA solution is only stable
when the per-edge aligned sample count exceeds a rank-dependent floor
(roughly $4k$ in our regularised implementation). An early iteration used
$k = 50$ and required $\approx 60$ aligned samples per edge---a floor that
realistic calibration windows could not always meet because inter-container
activity alignment is lossy. At $k = 15$ the floor drops to $\approx 25$
aligned samples; every edge observed in our testbed calibrates cleanly.
Section~\ref{sec:cal-output} reports that $k = 15$ retains more than 90\%
of the explained inter-container correlation.

%─────────────────────────────────────────────────────────────────────────────
\section{Mahalanobis edge energy}
\label{sec:mahal}
%─────────────────────────────────────────────────────────────────────────────

At runtime cycle $t$ the residual on edge $e = (u, v, \ell)$ is
$\mathbf{r}_e^{(t)} = F_u^{(e)} \tilde{\mathbf{x}}_u^{(t)} - F_v^{(e)}
\tilde{\mathbf{x}}_v^{(t+\ell)} \in \mathbb{R}^k$. Under the
calibrated-coupling null, $\mathbf{r}_e^{(t)}$ is zero-mean Gaussian in the
shared subspace. We estimate its covariance with Tikhonov regularisation:
\[
  \hat\Sigma_e = \frac{1}{T-1}\sum_{\tau=1}^{T}
  \mathbf{r}_e^{(\tau)}\bigl(\mathbf{r}_e^{(\tau)}\bigr)^{\!\top}
  + \epsilon_\Sigma I_k, \quad \epsilon_\Sigma = 10^{-6}.
\]
The \emph{edge energy} is the Mahalanobis distance~\cite{mahalanobis1936}:
\begin{equation}
  E_e(t) = \bigl(\mathbf{r}_e^{(t)}\bigr)^{\!\top} \hat\Sigma_e^{-1} \mathbf{r}_e^{(t)}.
  \label{eq:mahal}
\end{equation}
Under the null $E_e \sim \chi_k^2$; for $k{=}15$ the mean is $15$ and the
$99\%$ tail is $\approx 30.6$. The per-edge threshold is $\tau_e = \hat\mu_e
+ 4\hat\sigma_e$.

%─────────────────────────────────────────────────────────────────────────────
\section{Cellular sheaf Laplacian and Rayleigh quotient}
\label{sec:sheaf}
%─────────────────────────────────────────────────────────────────────────────

A cellular sheaf $\mathcal{F}$ on $G = (V, E)$ assigns a vector space
$\mathcal{F}(v) = \mathbb{R}^{d}$ to each vertex, a vector space
$\mathcal{F}(e) = \mathbb{R}^{k}$ to each edge, and a linear restriction
$\mathcal{F}_{v \triangleleft e} : \mathcal{F}(v) \to \mathcal{F}(e)$ to
every incident pair. The coboundary $\delta$ acts on an edge $(u, v)$ by
$(\delta \mathbf{x})_e = F_u^{(e)} \mathbf{x}_u - F_v^{(e)} \mathbf{x}_v$.
The \emph{sheaf Laplacian} is
\begin{equation}
  L_\mathcal{F} = \delta^{\!\top}\!\delta \in \mathbb{R}^{nd \times nd}.
  \label{eq:sheaf-laplacian}
\end{equation}
The diagonal block for vertex $v$ is $\sum_{e \colon v \triangleleft
e}(F_v^{(e)})^{\!\top} F_v^{(e)}$; the off-diagonal block for $(u,v)$
incident at edge $e$ is $-(F_u^{(e)})^{\!\top} F_v^{(e)}$. When every
restriction is the identity, $L_\mathcal{F}$ reduces to the ordinary graph
Laplacian~\cite{chung1997spectral}.

\paragraph{Intuition.}  Think of each container as a point in its own
74-dimensional configuration space, and each CCA restriction as the pair
of projection lenses that edge uses to see its two endpoints. If both
endpoints project onto the same point in the shared 15-dimensional
subspace, the edge is \emph{happy}. If they disagree, the edge has
\emph{tension}. The sheaf Laplacian sums tension across all edges. A low
Rayleigh quotient means the graph is at rest; a high one means the
coupling is being violated.

The global Rayleigh quotient is
\begin{equation}
  R(\mathbf{x}) = \frac{\mathbf{x}^{\!\top} L_\mathcal{F}\, \mathbf{x}}{\mathbf{x}^{\!\top} \mathbf{x}}
  = \frac{\sum_{e = (u,v) \in E} \|F_u^{(e)} \mathbf{x}_u - F_v^{(e)} \mathbf{x}_v\|_2^2}
         {\|\mathbf{x}\|_2^2}.
  \label{eq:rayleigh}
\end{equation}
We set $\tau_{\text{global}} = \hat\mu_R + 4\hat\sigma_R$ at calibration and
fire the global gate when $R > \tau_{\text{global}}$.

%─────────────────────────────────────────────────────────────────────────────
\section{Guarded exponential moving average}
\label{sec:ema}
%─────────────────────────────────────────────────────────────────────────────

A naive EMA adapts the threshold to current observations but has one fatal
mode: a slow-drip attack that never exceeds the instantaneous threshold
poisons the average into following the attack. Let $s^{(t)} \in \{0, 1\}$
indicate whether cycle $t$ was pristine (no behaviour bit fired, no edge
energy exceeded $\tau_e$, no alert in the ring buffer) and let $p^{(t)} =
\sum_{\tau = t-5}^{t} s^{(\tau)}$ be the pristine-streak length over the
last six cycles. The guarded update is
\begin{equation}
  \tau^{(t)} =
  \begin{cases}
    \alpha R^{(t)} + (1 - \alpha)\tau^{(t-1)} & \text{if } p^{(t)} = 6, \\[1mm]
    \tau^{(t-1)} & \text{otherwise}.
  \end{cases}
  \label{eq:guarded-ema}
\end{equation}
Any suspected anomaly freezes the threshold. We use $\alpha = 0.02$, giving
an effective half-life of roughly $34$ pristine cycles---slow enough to
track legitimate drift, fast enough to recover after a clean recovery
window.

%─────────────────────────────────────────────────────────────────────────────
\section{Evasion resistance: four theorems}
\label{sec:theorems}
%─────────────────────────────────────────────────────────────────────────────

The design of CausalTrace rests on four defensible claims. Each theorem
answers a specific question a reviewer or a panel member would raise. The
proofs are informal because each relies on a kernel-observable mechanism
rather than a closed-form reduction; the citations point to the kernel or
userspace file that realises the mechanism.

\begin{theorem}[Invariant necessity]\label{thm:invariant}
Let $\mathcal{A}_{\text{I}} \subseteq \mathcal{A}$ be the set of attack
classes that require the attacker to emit a specific syscall pattern $\pi$
to succeed (for example, an interactive reverse shell requires
$\mathrm{dup2}(\texttt{sockfd}, k)$ for some $k \in \{0,1,2\}$). For every
$\alpha \in \mathcal{A}_{\text{I}}$, a kernel hook observing $\pi$ achieves
recall $= 1$: the attacker cannot complete $\alpha$ without the hook
firing. Table~\ref{tab:invariants} maps each scenario to its invariant.
\end{theorem}

\begin{theorem}[Bigram-noise adversarial closure]\label{thm:bigram}
Let $N$ be the noise-syscall whitelist filtered before the bigram update
(Section~\ref{sec:signal}). An adversary attempting to dilute the
bigram-count vector during an attack faces a dichotomy:
\begin{enumerate}
  \item \textit{Emit syscalls in $N$.} These are filtered at the kernel and
  contribute nothing to the sketch. They cost CPU and nothing else.
  \item \textit{Emit syscalls outside $N$.} Each such syscall produces a
  bigram $(\mathrm{prev}, s)$ whose probability under the calibrated model
  is low (it was not seen during calibration). The resulting
  $\tilde{\mathbf{x}}^{(t)}$ moves farther from the calibrated mean, the
  Mahalanobis edge energy~\eqref{eq:mahal} rises, and the sheaf Rayleigh
  quotient rises with it.
\end{enumerate}
The attacker has no noise-channel at the bigram layer: every evasion
attempt triggers the statistics.
\end{theorem}

\begin{theorem}[Two-layer completeness]\label{thm:completeness}
Partition the attack space as $\mathcal{A} = \mathcal{A}_{\text{I}} \cup
\mathcal{A}_{\text{U}}$, where $\mathcal{A}_{\text{U}}$ is the set of novel
attacks without a pre-enumerated kernel invariant. Theorem~\ref{thm:invariant}
gives Tier~1 recall $= 1$ on $\mathcal{A}_{\text{I}}$. Tier~3 achieves
recall $\geq 1 - \beta$ on $\mathcal{A}_{\text{U}}$, where $\beta$ is
controlled by the per-edge false-positive rate $\Pr[E_e > \tau_e \mid
\text{null}]$ and the $4\sigma$ gate in Sections~\ref{sec:mahal}
and~\ref{sec:sheaf}. The composition covers $\mathcal{A}$ with detection
rate bounded below by the invariant coverage on the known set and the
sheaf's statistical power on the novel set.
\end{theorem}

\begin{theorem}[Trust-promotion soundness]\label{thm:trust}
Let $B = 5120$ bytes and $\Delta = 1$~second be the L4-stability gate for
trust promotion (Chapter~\ref{Chapter4}, Section~\ref{sec:trust-model}). A
client IP reaches \textsc{Calibrated} only after a single TCP flow
accumulates at least $B$ bytes over at least $\Delta$ seconds. An attacker
whose entire session budget is under $B$ bytes cannot promote, and
therefore cannot bypass Case~A enforcement; a burst attacker who fires and
disconnects before $\Delta$ elapses is always routed through the
untrusted path.
\end{theorem}

\begin{remark}[Unforgeable attribution]\label{rem:attribution}
The per-TID client-IP map is populated by the \texttt{inet\_csk\_accept}
kretprobe using kernel-owned \texttt{struct sock} fields. No userspace
write path exists to this map; an attacker cannot spoof a client-IP entry
without kernel-level privilege escalation, which is outside the threat
model of Chapter~\ref{Chapter4}.
\end{remark}

\paragraph{What we do \emph{not} prove.}  The sheaf detector's
$\beta$ in Theorem~\ref{thm:completeness} is an empirical quantity---we do
not claim a closed-form upper bound. We do not claim that
Theorem~\ref{thm:invariant} holds for attack classes that a future kernel
may allow through new syscalls not yet in $\Sigma$. Theorem~\ref{thm:trust}
assumes the attacker cannot fabricate TCP traffic on the defender's behalf,
which holds as long as the defender's veth is the only way into the
container.

%─────────────────────────────────────────────────────────────────────────────
\section{Math-to-code mapping}
\label{sec:math-map}
%─────────────────────────────────────────────────────────────────────────────

\begin{table}[h]
\centering
\caption[Math-to-code mapping]{Where each mathematical object is realised
in the codebase.}
\label{tab:math-map}
\small
\begin{tabular}{lll}
\toprule
Object / equation & Kernel side & Userspace side \\
\midrule
Bigram CMS, $4 \times 128$                       & \texttt{kernel/causaltrace\_bcc.c:613} & \texttt{tier3/signal\_extractor.py} \\
Signal extraction $\phi$, $d{=}74$               & --- & \texttt{tier3/signal\_extractor.py} \\
Per-container whitening                          & --- & \texttt{tier3/whitener.py} \\
CCA restriction maps~\eqref{eq:cca}              & --- & \texttt{tier3/calibrate.py} \\
Mahalanobis energy~\eqref{eq:mahal}              & --- & \texttt{tier3/sheaf\_detector.py} \\
Sheaf Laplacian~\eqref{eq:sheaf-laplacian}       & --- & \texttt{tier3/sheaf\_detector.py} \\
Rayleigh quotient~\eqref{eq:rayleigh}            & --- & \texttt{tier3/sheaf\_detector.py} \\
Guarded EMA~\eqref{eq:guarded-ema}               & --- & \texttt{tier3/ema\_buffer.py} \\
Invariant hooks (Thm.~\ref{thm:invariant})       & \texttt{kernel/causaltrace\_bcc.c} & --- \\
Trust FSM (Thm.~\ref{thm:trust})                 & \texttt{kernel/causaltrace\_bcc.c:459} & \texttt{tier3/trust\_promoter.py} \\
\bottomrule
\end{tabular}
\end{table}

The per-scenario invariant map (Table~\ref{tab:invariants}) is deferred to
Chapter~\ref{Chapter4} where the full attack catalogue is introduced.
